\chapter[Applications to projective algebraic schemes]
{Applications\\ to projective algebraic schemes} \label{XII}

\renewcommand{\theequation}{\arabic{equation}}

\section[Projective duality theorem and finiteness theorem]
{Projective duality theorem and finiteness theorem\protect\sfootnotemark}\label{XII.1}

The\pageoriginale following theorem,\sfootnotetext{The present expos\'e, written in January~1963, is considerably more detailed than was the oral presentation in June~1962.} essentially contained in FAC\sfootnote{J.-P.~Serre, \sisi{Faisceaux alg\'ebriques Coh\'erents, Ann. of Math. t.61 (1955) 197-278}{{\og Faisceaux alg\'ebriques coh\'erents\fg}, \emph{Ann. of Math.} \textbf{61} (1955), p\ptbl197-278}.} (except that, at the time, Serre did not have at his disposal the language of Ext for sheaves of modules\nde{Readers interested in the history of mathematics will profitably consult the letter that Grothendieck wrote to Serre on 15 December 1955 and Serre's reply of 22 December of the same year; see \emph{Correspondance Grothendieck-Serre}, edited by Pierre Colmez and Jean-Pierre Serre, Documents Math\'ematiques, vol.~2, Soci\'et\'e Math\'ematique de France, Paris, 2001.}), is the global analogue of the local duality theorem (Expos\'e \Ref{IV}), which was modeled on it.

\begin{theoreme} \label{XII.1.1}
Let $k$ be a field, $X=\PP_k^r$ \sisi{the standard projective scheme}{projective space} of dimension $r$ over $k$, and $F$~a variable coherent module on $X$; then there is an isomorphism of $\partial$-functors in~$F$:
\steco
\begin{equation} \label{eq:XII.1} {}\H^i(X, F)' \isomfrom \Ext^{r-i}(X;F, \omx^r),
\end{equation}
where we set
\steco
\begin{equation} \label{eq:XII.2} {}\Omega_{X/k}^r= \Oo_{\PP_k^r}(-r-1).
\end{equation}
\end{theoreme}

\begin{remarquestar}
This module is of course also the module of relative differentials of degree $r$ of $X$ over $k$. In this form, the theorem remains valid if $X$ is a proper smooth scheme over $k$ (for the projective case, see A. Grothendieck, \og \sisi{Th\'eor\`emes de dualit\'e pour les faisceaux alg\'ebriques coh\'erents}{\emph{Th\'eor\`emes de dualit\'e pour les faisceaux alg\'ebriques coh\'erents}}\fg, S\'eminaire Bourbaki, May 1957)\sfootnote{For a more general duality theorem, \cf the Hartshorne seminar cited at the end of \Exp \Ref{IV}.}. When $F$ is locally free, one recovers Serre's duality theorem $\H^i(X, F)' \isomto \H^{r-i}(\SheafHom_{\OX}(F, \omx^r))$. Theorem \Ref{XII.1.1} (which also recovers the case of a projective $X$ over $k$, as in {\loccit}) will suffice for our purposes.
\end{remarquestar} The homomorphism\pageoriginale (\Ref{eq:XII.1}) is deduced from the Yoneda pairing
\steco
\begin{equation} \label{eq:XII.3} {}\H^i(X, F) \times \Ext^{r-i}(X;F, \omx^r) \to \H^r(X, \omx^r),
\end{equation}
and from a well-known isomorphism (\cf FAC, or \EGA III 2.1.12):
\steco
\begin{equation} \label{eq:XII.4} {}\H^r(X, \omx^r) = \H^r(\PP_k^r, \Oo_{\PP_k^r}(-r-1)) \isomto k.
\end{equation}
To show that (\Ref{eq:XII.1}) is an isomorphism, one proceeds as in the case of the local duality theorem, noting that $\H^r(X, F)$, as a functor of $F$, is right exact (since $\H^n(X, F)=0$ for $n>r$), and that every coherent module is isomorphic to a quotient of a direct sum of modules of the form $\Oo(-m)$, with $m$ large. By descending induction on $i$, this reduces us to checking the assertion for a sheaf of the form $\Oo(-m)$, where it is contained in the well-known explicit computations (FAC, or \EGA III 2.1.12). We may moreover assume $-m \leq -r-1$, in which case $\H^i(X, \Oo(-m))=0$ for $i\neq r$.

\begin{corollaire} \label{XII.1.2} For a given coherent $F$, and for $m$ sufficiently large, there is a canonical isomorphism
\steco
\begin{equation} \label{eq:XII.5} {}\H^i(X, F(-m))'\isomto\H^0(X, \SheafExt_{\OX}^{r-i}(F, \omx^r)(m))
\end{equation}
(where $'$ denotes the vector-space dual).
\end{corollaire}
Indeed, on \sisi{a projective scheme}{projective space} $X$, for every pair of coherent sheaves $F,G$ and for $n$ sufficiently large, there is a canonical isomorphism:
\steco
\begin{equation} \label{eq:XII.6} {}\Ext^n(X;F(-m), G)\simeq \Ext^n(X;F, G(m)) \isomto \H^0(X, \SheafExt_{\OX}^{n}(F, G)(m)),
\end{equation}
(the isomorphism between the first two terms being trivially valid for every $m$), as follows from the global $\Ext$ spectral sequence
\[
\H^p(X, \SheafExt^q_{\OX}(F, G(m))) \To \Ext^{\boule}(X;F, G(m)),
\]
which\pageoriginale degenerates for $m$ sufficiently large, by virtue of the fact that
\[
\SheafExt^q_{\OX}(F, G(m)) \simeq \SheafExt^q_{\OX}(F, G)(m),
\]
and that the $\SheafExt^q_{\OX}(F, G)$ are coherent sheaves. Thus (\Ref{eq:XII.5}) follows from (\Ref{eq:XII.6}) and (\Ref{eq:XII.1}).

\begin{corollaire} \label{XII.1.3}
\sisi{Equivalent conditions for given $i,F$}{For given $i,F$, the following conditions are equivalent}:
\begin{enumeratei}
\item
$\H^i(X, F(-m))=0$ for $m$ large.
\item[\textup{(i bis)}] $\H^i(X, F(\cdot))= \sisi{\coprod}{\bigoplus}_{m\in\ZZ}\H^i(X, F(-m))$ is a finite-type $S$-module, where $S=k[t_o, \ldots, t_r]$.
\item
$\SheafExt^{r-i}_{\OX}(F, \omx^r)=0$.
\item[\textup{(ii bis)}] $\SheafExt^{r-i}_{\OX}(F, \OX)=0$.
\item
$\H^i_x(F_x)=0$ for every closed point $x$ of $X$.
\item
$\H^{i+1}_x(\widetilde{F}_x)=0$ for every closed point $x$ of the punctured affine cone $\widetilde{X}=\Spec(S)-\Spec(k)$ over $X$, where $\widetilde{F}$ denotes the inverse image of $F$ under the canonical morphism $\widetilde{X}\to X$.
\end{enumeratei}
\end{corollaire}

\begin{proof}
(i) \SSI (i bis), since the submodule of $\H^i(X, F(\cdot))$ that is the sum of the homogeneous components of degree $\geq \nu$ is of finite type over $S$ (in fact, for $i\neq 0$, it is even of finite type over $k$), (\cf FAC or \EGA III 2.2.1 and 2.3.2).

(i) \SSI (ii) by Corollary \Ref{XII.1.2}.

(ii) \SSI (ii bis), because $\Omega_{X/k}^r$ is locally isomorphic to $\OX$.

(ii bis) \SSI (iii) by the local duality theorem for $\OXx$ (which is a regular\pageoriginale local ring of dimension $r$), according to which the \og dual\fg of $\SheafExt^{r-i}_{\OX}(F, \OX)_x$ is identified with $\H^i_x(F_x)$ (\hyperref[V.2.1]{V 2.1)}.

(ii bis) is equivalent to the analogous relation
\[
\SheafExt^{r-i}_{\Oo_{\widetilde{X}}}(\widetilde{F}, \Oo_{\widetilde{X}})=0
\]
(because $\widetilde{X}\to X$ is faithfully flat, so the inverse image of $\SheafExt^{r-i}_{\OX}(F, \OX)$ is isomorphic to $\SheafExt^{r-i}_{\Oo_{\widetilde{X}}}(\widetilde{F}, \Oo_{\widetilde{X}})$), and this last relation is equivalent to (iv) by the local duality theorem for the local ring $\OXx$, which is regular of dimension $r+1$.
\skipqed
\end{proof}

In particular, applying this to all $i\leq n$, we find:

\begin{corollaire} \label{XII.1.4}
Equivalent conditions for given $n,F$:
\begin{enumeratei}
\item
$\H^i(X, F(-m))=0$ for $i\leq n$ and $m$ large.
\item[\textup{(i bis)}] $\H^i(X, F(\cdot))$ is a finite-type $S$-module for $i\leq n$.
\item
$\prof(F_x)>n$ for every closed point $x$ of $X$.
\item
$\prof(\widetilde{F}_x)>n+1$ for every closed point $x$ of $\widetilde{X}$.
\end{enumeratei}
\end{corollaire}

The point of Corollaries \Ref{XII.1.3} and \Ref{XII.1.4} is to express a global condition (i) or~(i~bis) in terms of local conditions, \sisi{\ignorespaces}{namely} the vanishing of local invariants such as $\H_x^i(X, F_x)$ or $\H_x^i(X, \widetilde{F}_x)$, or an inequality involving depth. In this form, these results remain trivially valid for an arbitrary projective scheme $X$, and a very ample invertible sheaf $\OX(1)$ on $X$, as one sees by means of a suitable projective immersion $X\to\PP_k^r$. (Of course, conditions \Ref{XII.1.3}~(ii) and \Ref{XII.1.3}~(ii bis) are no longer equivalent to the others in this general case, unless, for example, one assumes that $X$ is regular.) One can moreover generalize to the case of a projective morphism $X\to S$ as follows:

\begin{proposition} \label{XII.1.5}
Let\pageoriginaled $f:X\to S$ be a projective morphism, with $S$ noetherian, let $\OX(1)$ be an invertible module on $X$ that is very ample relative to $S$, let $F$ be a coherent module on $X$, flat over $S$, let $s$ be a point of $S$, let $X_s$ be the fiber of $X$ at $s$ (regarded as a projective scheme over $k(s)$), let $F_s$ be the sheaf induced on $X_s$ by $F$, and finally let $i$ be an integer. Suppose that, for every closed point $x$ of $X_s$, one has $\H^i_x(F_{s, x})=0$ (for example, $\prof(F_{s, x})>i$). Then there exists an open neighborhood $U$ of $s$ such that the same condition holds for every $s'\in U$. Moreover, for such a $U$, one has
\[
\R^if_\ast(F(-m))= 0 \; \text{for}\; m \; \text{large, }
\]
and if $\Ss$ is a quasi-coherent graded algebra on $S$, generated by $\Ss^1$, and defining~$X$ together with $\OX(1)$ as $X=\Proj(\Ss)$, $\OX(1)=\Proj(\Ss(1))$, then the $S$-module
\[
\R^if_\ast(F(\cdot)) = \sisi{\coprod}{\tbigoplus}_{m\in\ZZ}\R^if_\ast(F(m))
\]
is of finite type over $U$.
\end{proposition}

Embed $X$ in an $X'=\PP_S^r$ in such a way that $\OX(1)$ is induced by $\Oo_{\sisi{P}{{X'}}}(1)$ (this is possible provided that $S$ is replaced by an affine neighborhood of $s$). Put\refstepcounter{toto}\label{P5}\sfootnote{The first part of \Ref{XII.1.5} can also be obtained immediately by applying the purely local statement \EGA IV 12.3.4 to the preceding $\underline{E}^j$, thereby bypassing most of the proof that follows.} for every integer $j$, and every $t\in S$:
\steco
\begin{equation} \label{eq:XII.7} {}\underline{E}^j(t)= \SheafExt^j_{\Oo_{X'_t}} (F_{t'}, \Oo_{X'_t}(-r-1))
\end{equation}
Thus $\underline{E}^j(t)$ is a coherent module on $X_t$; I claim that, as $t$ varies, the family of these modules is \og constructible\fg in the following sense: for every $t\in S$ there exists a nonempty open subset $V$ of the closure $\overline{t}$, endowed with the induced reduced structure, and a coherent module $\underline{E}^j(V)$ on $X_V=X\times_SV$, \emph{flat} over $V$, such that, for every $t'\in V$, $\underline{E}^j(t)$ is isomorphic to the module induced by $\underline{E}^j(V)$ on $X_t$. To verify this assertion, putting $Z=\overline{t}$ with its induced structure, we consider the coherent modules\pageoriginale
\[
\underline{E}^j(Z)= \SheafExt^j_{\Oo_{X'_Z}} (F_Z, \Oo_{X'_Z}(-r-1))
\]
(where the subscript $Z$ still means that one takes the induced object \sisi{over}{over} $Z$), and we take for $V$ a nonempty open subset of $Z$ such that the modules $\underline{E}^j(Z)$ are flat \sisi{over}{over} $V$: this is possible, since one verifies immediately that $\underline{E}^j(Z)=0$ for $j$ outside the interval $[0, r]$, and one can apply \SGA 1 IV 6.11. We then take $\underline{E}^j(V)=\underline{E}^j(Z)|_{X_V}$, and one readily verifies that it has the required property.

It follows from the preceding remark that there exists a finite partition of $S$ formed by sets $V_\alpha$ of the form $V=V(t)$ as above (noetherian induction), and, applying Serre's theorem \EGA III 2.2.1 to the $\underline{E}^j(V_i)$, one sees that there exists an integer $m_0$ such that
\[
\R^if_{V_\alpha\ast}(\underline{E}^j(V_\alpha))=0\textup{ for } i\neq 0, m\geq m_0, \textup{ for }j,
\]
from which it follows, using the flatness of $\underline{E}^j(V_\alpha)$ over $V_\alpha$ and straightforward K\"unneth-type relations (\cf \EGA III par\ptbl7), that
\[
\H^i(X_t, \underline{E}^j(t)(m))=0\textup{ for } i\neq 0, m\geq m_0, \textup{ for }j,
\]
for every $t\in V_\alpha$, and hence for every $t$, since the $ V_\alpha$ cover $S$. From this and the spectral sequence for global Ext it follows, by virtue of \Ref{XII.1.1}, as in the proof of \Ref{XII.1.2}, that there is an isomorphism
\steco
\begin{equation} \label{eq:XII.8} {}\H^i(X_t, F_t(-m))' \isomfrom \H^0(X_t, \underline{E}^{r-i}(t)(m))\textup{ for } m\geq m_0,
\end{equation}
for every integer $i$, and every $t\in S$.

Now use the hypothesis on $F_s$, which can be written
\steco
\begin{equation} \label{eq:XII.9} {}\underline{E}^{r-i}(s)=0,
\end{equation}
and which, by virtue of (\Ref{eq:XII.8}), is equivalent to
\steco
\begin{equation} \label{eq:XII.10} {}\H^i(X_s, F_s(-m))=0 \; \text{for}\; m\geq m_0.
\end{equation}
Since $F$, and hence $F(-m)$, is flat over $S$, it follows from
the K\"unneth-type relations already invoked that (for a given $m$) the
same relation (\Ref{eq:XII.10}) holds when\pageoriginale $s$ is replaced by a
$t$ near $s$, in particular by every generization $t$ of
$s$. By virtue of (\Ref{eq:XII.8}), we shall therefore have, for such a
generization, \steco
\begin{equation} \label{eq:XII.11} {}\underline{E}^{r-i}(t)=0,
\end{equation}
and the set of $t\in s$ for which this relation holds is plainly constructible (since it induces an open subset on each $V_\alpha$); as it contains the generizations of $s$, it contains an open neighborhood $U$ of $s$. This proves the first assertion of \Ref{XII.1.5}. Moreover, for $t\in U$, it follows from (\Ref{eq:XII.11}) and (\Ref{eq:XII.8}) that
\steco
\begin{equation} \label{eq:XII.12} {}\H^i(X_t, F_t(-m)=0 \; \text{for}\; m\geq m_0, \; t\in U,
\end{equation}
which, by virtue of the K\"unneth-type relations, implies (and in fact is much stronger than)
\steco
\begin{equation} \label{eq:XII.13} {}\R^if_\ast(F(-m))=0 \; \text{on}\; U, \; \text{for} \; m\geq m_0.
\end{equation}
This proves the second assertion of \Ref{XII.1.5}. Finally, the last assertion follows immediately by proceeding as at the beginning of the proof of \Ref{XII.1.3}.

\refstepcounter{subsection}\label{XII.1.6}
\begin{enonce*}[remark]{Remark 1.6\sfootnotemark} \sfootnotetext{This remark is made more precise by footnote \sisi{5.}{\pageref{P5}.}} The proof simplifies considerably (eliminating all consideration of constructibility) when one assumes from the outset that the hypothesis made for $s\in S$ holds for every $s'\in S$. In fact, when one assumes that $F_s$ has depth $>i$ at the closed points of $X$, there is a general statement, \emph{local in nature on} $X$, which says that the same condition holds for all the $X_t$, provided that $X$ is replaced by a suitable open neighborhood of the fiber $X_s$ (in other words, a certain subset of $X$, defined by conditions on the modules induced by $F$ on the fibers, is open, \cf \EGA IV). Since $f$ is proper here, this neighborhood can therefore be taken to be of the form $f^{-1}(U)$, which recovers the first assertion of \Ref{XII.1.5} without laborious d\'evissage. In this general case, one can still prove by the method of \loccit that the first assertion of \Ref{XII.1.5} (proved here globally, using the fact that $X$ is projective over $S$) again follows from a purely local statement on $X$ (which the reader may formulate explicitly if useful).
\end{enonce*}

\section[Lefschetz theory for a projective morphism: comparison]
{Lefschetz theory for a projective morphism: Grauert comparison theorem} \label{XII.2}

\setcounter{equation}{13}
It is\pageoriginale the following theorem:

\begin{theoreme} \label{XII.2.1}
Let $f:X\to S$ be a projective morphism, with $S$ Noetherian, let $\OX(1)$ be an invertible module on $X$, relatively ample over $S$, let $Y$ be the prescheme of zeros of a section $t$ of $\OX(1)$, let $J$ be the ideal defining $Y$, let $X_n$ be the closed subprescheme of $X$ defined by $J^{n+1}$, let $\widehat{X}$ be the formal completion of $X$ along $Y$, let $\widehat{f}:\widehat{X}\to S$ be the composite morphism $\widehat{X}\to X\to S$, and let $F$ be a coherent module on $X$, flat over $S$. Assume in addition that, for every $s\in S$, the module $F_s$ induced on the fiber $X_s$ has depth $>n$ at the points of that fiber, and that $t$ is $F$-regular. Under these conditions:
\begin{enumeratei}
\item
The canonical homomorphism
\[
\R^if_\ast(F) \to \R^i\widehat{f}_\ast(\widehat{F})
\]
is an isomorphism for $i<n$, and a monomorphism for $i=n$.
\item
The canonical homomorphism
\[
\R^i\widehat{f}_\ast(\widehat{F}) \to \varprojlim_m \R^if_\ast(F_m)
\]
is an isomorphism for $i\leq n$.
\end{enumeratei}
\end{theoreme}

\begin{proof}
We immediately reduce to the case where $S$ is affine, and then to proving the

\begin{corollaire} \label{XII.2.2}
Under the conditions of \Ref{XII.2.1}, assume in addition that $S$ is affine. Then:
\begin{enumeratei}
\item
The canonical homomorphism
\[
\H^i(X, F) \to \H^i(\widehat{X}, \widehat{F})
\]
is\pageoriginale an isomorphism for $i<n$, and a monomorphism for $i=n$.
\item
The canonical homomorphism
\[
\H^i(\widehat{X}, \widehat{F}) \to \varprojlim_m \H^i(X_m, F_m)
\]
is an isomorphism for $i\leq n$.
\end{enumeratei}
\end{corollaire}

After replacing $\OX(1)$ by a tensor power, and $t$ by a power of~$t$, we may assume that $\OX(1)$ is very ample relative to $S$. On the other hand, since $t$, and hence $t^m$, is $F$-regular, multiplication by $t^m$, regarded as a homomorphism from $F(-m)$ to $F$, is injective; therefore, for every $m\geq 0$, there is an exact sequence:
\steco
\begin{equation} \label{eq:XII.14}
0\to F(-m) \lto{t^m} F \to F_m \to0,
\end{equation}
and hence an exact cohomology sequence
\[
\H^i(X, F(-m)) \to \H^i(X, F) \to \H^i(X, F_m) \to \H^{i+1}(X, F(-m)).
\]
But by \Ref{XII.1.5}, we have $\H^i(X, F(-m))=0$ for $i\leq n$ and $m$ sufficiently large, which proves the

\begin{lemme} \label{XII.2.3}
For $m$ large, the canonical homomorphism
\[
\H^i(X, F) \to \H^i(X, F_m)
\]
is bijective if $i<n$, and injective if $i=n$.
\end{lemme}

\enlargethispage{-2\baselineskip}%
This shows that, for $i<n$, the inverse system $(\H^i(X_m, F_m))_{m\geq 0}$ is essentially constant, and \afortiori satisfies the Mittag-Leffler condition; hence (\sisi{given that}{in view of} $\widehat{F}= \varprojlim F_m$), we obtain (ii) from $\EGA 0_{\textup{III}}$ 13.3. On the other hand, (i) follows trivially from this, taking \Ref{XII.2.3} into account.

\refstepcounter{subsection}\label{XII.2.4}
\begin{enonce*}{Corollary 2.4\ndemark}\ndetext{See Corollary I.1.4 of the article by Mme Raynaud (Raynaud~M., {\og Th\'eor\`emes de Lefschetz en cohomologie des faisceaux coh\'erents et en cohomologie \'etale. Application au groupe fondamental\fg}, \emph{Ann. Sci. \'Ec. Norm. Sup. (4)} \textbf{7} (1974), p\ptbl 29--52).}
Let\pageoriginale $f:X\to S$ be a projective flat morphism, with $S$ locally Noetherian, let $\OX(1)$ be an invertible module on $X$, relatively ample over $S$, let $t$ be a section of this module that is $\OX$-regular, let $Y$ be the closed subprescheme of zeros of $t$, and let $\widehat{X}$ be the formal completion of $X$ along $Y$. Assume that, for every $s\in S$, $X_s$ has depth $\geq 1$ (\resp depth $\geq 2$) at its closed points. Then, for every open neighborhood $U$ of $Y$, the functor
\[
F \mto \widehat{F}
\]
from the category of locally free coherent modules on $U$ to the category of locally free coherent modules on $\widehat{X}$ is faithful (\resp fully faithful, \ie the Lefschetz condition ($\Lef$) of \Ref{X}~\Ref{X.2} holds).
\end{enonce*}

For two locally free Modules $F$ and $G$ on $U$, introduce the Module
\[
H=\SheafHom_{\Oo_U}(F, G)
\]
we are reduced to proving that the canonical homomorphism
\steco
\begin{equation} \label{eq:XII.15} {}\H^0(U, H) \to \H^0(\widehat{U}, \widehat{H})
\end{equation}
is injective (\resp bijective). Now the Modules $H_t$ have depth $\geq 1$ (\resp $\geq 2$) at the closed points of $X_t$, and hence we can apply \Ref{XII.2.1}, which gives the conclusion \Ref{XII.2.4} in the case where $U=X$. In the case of an arbitrary $U$, note that the question is local on $S$, so we may suppose that $S$ is affine. Then every coherent Module on $X$ is a quotient of a coherent locally free Module (since $\OX(1)$ is an invertible \sisi{absolutely}{relatively} ample Module on $X$). Since the dual Module $H'=\SheafHom(H, \Oo_U)$ extends to a coherent Module on $X$, which is therefore isomorphic to a cokernel of a homomorphism of locally free Modules on $X$, it follows by transposition that one can find a homomorphism
\[
u':{L'}^0 \to {L'}^1
\]
of\pageoriginale locally free Modules on $X$, inducing a homomorphism
\[
u:L^0 \to L^1
\]
of locally free Modules on $U$, such that there is an exact sequence
\[
0\to H \to L^0 \lto{u} L^1.
\]
Using the five lemma (which becomes the three lemma), and the left exactness of the functor $H^0$, we are reduced to proving that (\Ref{eq:XII.15}) is injective (\resp bijective) when $H$ is replaced there by $L^0$, $L^1$, which brings us back to the case where $H$ is induced by a locally free Module $H'$ on $X$. Moreover, in the non-parenthetical case this reduction is even unnecessary, since the kernel of (\Ref{eq:XII.15}) consists in any event of the sections of $H$ on~$U$ that vanish in some suitable open neighborhood $V$ of $Y$; but the restriction homomorphism $\H^0(U, H)\to\H^0(V, H)$ is injective, since $H$ has depth $\geq 1$ at the points of every closed subset $Z$ of $X$ not meeting $Y$ (\cf the lemma below). In the parenthetical case, we are reduced to proving that
\[
H^0(X, H')\to\H^0(U, H')
\]
is bijective, which follows from the fact that $H'$ has depth $\geq 2$ at all the points of a closed subset $Z=X- U$ of $X$ not meeting $Y$. It therefore remains only to prove the

\begin{lemme} \label{XII.2.5} Let $F$ be a coherent Module on $X$, flat over $S$, such that for every $s\in S$, $F_s$ has depth $\geq n$ at every closed point of $X_s$. Then, for every closed subset $Z$ of $X$ not meeting $Y$, $F$ has depth $\geq n$ at all the points of $Z$.
\end{lemme}

Indeed, for every $x\in X$, setting $s=f(x)$, we have
\steco
\begin{equation} \label{eq:XII.16} {}\prof(F_x)\geq \prof(F_{s, x}),
\end{equation}
as can be seen by lifting arbitrarily a maximal $F_{s, x}$-regular sequence of elements of\pageoriginale $\rr(\Oo_{X_s, x})$, which gives an $F_x$-regular sequence by \SGA 1 IV 5.7. Now if $x$ belongs to a $Z$ as in Lemma~\Ref{XII.2.5}, then $x$ is necessarily closed in $X_s$; in other words, $Z$ is \emph{finite} over $S$. Indeed, $Z$ (endowed with a structure induced from $X$) is projective over $S$ as a closed subprescheme of $X$, which is projective, and $Z$ is affine over $S$ as a closed subprescheme of $X- Y$, which is affine.
\skipqed
\end{proof}

\begin{remarque} \label{XII.2.6}
Suppose that for every $s\in S$, the section $t_s$ of $\Oo_{X_s}(1)$ induced by $t$ is $\Oo_{X_s}$-regular (which implies by \SGA 1 IV 5.7 that $t$ is $\OX$-regular). Then the hypotheses made are stable under extension of the base $S'\to S$ ($S'$ locally noetherian). Thus the conclusion remains valid after every base change.
\end{remarque}

\section[Lefschetz theory for a projective morphism: existence]{Lefschetz theory for a projective morphism: existence theorem}\label{XII.3}%

\setcounter{equation}{16}%
\refstepcounter{subsection}\label{XII.3.1}%
\begin{enonce*}{Theorem 3.1\ndemark}
\ndetext{For a version without the flatness hypothesis, see (Raynaud~M., {\og Th\'eor\`emes de Lefschetz en cohomologie des faisceaux coh\'erents et en cohomologie \'etale. Application au groupe fondamental\fg}, \emph{Ann. Sci. \'Ec. Norm. Sup. (4)} \textbf{7} (1974), p\ptbl 29--52, Theorem II.3.3).}
Let $f:X\to S$ be a projective morphism, with $S$ noetherian, let $\OX(1)$ be an invertible Module on $X$ that is ample relative to $S$, let $X_0$ be the closed subprescheme of zeros of a section $t$ of $\OX(1)$, let $\widehat{X}$ be the formal completion of $X$ along $X_0$, let $\formelF$ be a coherent Module on $\widehat{X}$, and let $F_0$ be the Module that it induces on $X_0$. Assume moreover:
\begin{enumeratea}
\item\label{XII.3.1a}
$\formelF$ is flat over $S$.
\item\label{XII.3.1b}
For every $s\in S$, the section $t_s$ induced by $t$ on $X_s$ is $\formelF_s$-regular (which implies that $F_0$ is also flat over $S$, \cf \textup{\SGA 1 IV 5.7}).
\item\label{XII.3.1c}
For every $s\in S$, $F_{0s}$ has depth $\geq 2$ at the closed points of $X_{0s}$.
\end{enumeratea}
Assume moreover that $S$ admits an ample invertible sheaf. Under these conditions, there exists a coherent Module $F$ on $X$, and an isomorphism from its formal completion $\widehat{F}$ to $\formelF$.
\end{enonce*}

This statement will follow from the following one:

\begin{corollaire} \label{XII.3.2}
Under conditions \Ref{XII.3.1a}), \Ref{XII.3.1b}), \Ref{XII.3.1c}) above, the following hold:
\begin{enumeratei}
\item
The\pageoriginale Module $\widehat{f}_\ast(\formelF)$ on $S$ is coherent; hence, for every $n$, the Module $\widehat{f}_\ast(\formelF(n))$ on~$S$ is coherent.
\item
For $n$ large, the canonical homomorphism $\widehat{f}^\ast\widehat{f}_\ast(\formelF(n))\to \formelF(n)$ is surjective.
\end{enumeratei}
\end{corollaire}

Assume the corollary for the moment, and let us prove \Ref{XII.3.1}. By virtue of the last hypothesis made in \Ref{XII.3.1}, we may reduce to the case $X=\PP^r_S$, after replacing $\OX(1)$ and $t$ by a suitable power. I claim that we may further suppose that $t_s\neq 0$ for every~$s$. Otherwise, indeed, $\formelF_s=0$ by b), or equivalently, by Nakayama, $F_{0s}=0$, \ie $s$ does not belong to the image of $\Supp F_0$ under the morphism $f_0:X_0\to S$ induced by $f$. Now this image $S'$ is open by virtue of a), b), since $F_0$ is flat over $S$, and it is clear that it suffices to prove the conclusion of \Ref{XII.3.1} in the situation obtained by restricting over $S'$, since the coherent Module $F'$ on $X|_{S'}$ thereby obtained will be the restriction of a coherent Module $F$ on $X$, which will have the required property. We may therefore suppose that, in addition to hypotheses a), b), c), the following hypotheses also hold:
\begin{enumerate}
\item[\textup{a$'$)}] $\OX$ is flat over $S$.
\item[\textup{b$'$)}] For every $s\in S$, the section $t_s$ is $\Oo_{X_s}$-regular.
\item[\textup{c$'$)}] For every $s\in S$, $\Oo_{X_{0s}}$ has depth $\geq 2$ at the closed points of $X_{0s}$. (It~suffices to choose $X=\PP_S^r$ with $r\geq 3$, which is permissible.)
\end{enumerate}
Now \Ref{XII.3.2} implies that one can find an epimorphism
\steco
\begin{equation} \label{eq:XII.17}
\widehat{L} \to \formelF \to 0,
\end{equation}
where $L$ is a Module on $X$ of the form $f_\ast(G)(-n)$, with $G$ a coherent locally free Module on $S$: indeed, for $n$ large, it suffices to express the coherent Module $\widehat{f}_\ast(\formelF)$ on $S$ as a quotient of such a $G$. On the other hand, hypotheses a), b), c) on $f$, $t$ imply that $\widehat{L}$ satisfies the same conditions a), b), c) as $\formelF$. It follows easily that the same is true of the kernel of (\Ref{eq:XII.17}), to which the same argument may therefore be applied, so that $\formelF$ is represented as the cokernel of a homomorphism
\steco
\begin{equation} \label{eq:XII.18}
\widehat{L'} \to \widehat{L},
\end{equation}
where $L$, $L'$ are locally free Modules on $X$. Now, by virtue of a$'$), the second part of c$'$), and either \Ref{XII.2.1} or \Ref{XII.2.4}, the homomorphism (\Ref{eq:XII.18}) comes from a homomorphism $L'\to L$ of Modules on $X$. It now suffices to take for $F$ the cokernel of $L'\to L$, and we are done.

It remains to prove \Ref{XII.3.2}. This had been done in the
seminar by a somewhat laborious expedient, consisting in interpreting
everything in terms of cohomology on the punctured projective cone
of $X$ relative to $S$, in order to reduce to
Theorem \Ref{IX.2.1}. A more direct and more
satisfactory method (although substantially identical) now seems to me
to be the following. It consists in noting that in \Ref{IX},
\numero \Ref{IX.2}\refstepcounter{toto}\label{XII.14} (and with the
notation of that expos\'e), the hypothesis that the morphism
$f:\Xx\to\Xx'$ is adic is not used anywhere in the
proof of \Ref{IX.2.1}, via $\EGA 0_{\textup{III}}$ 13.7.7;
it suffices instead to suppose that $\Xx$ is also adic,
and to choose two ideals of definition $\Jj$ for $\Xx'$, $\Ii$
for $\Xx$, such that $\Jj\Oo_\Xx\subset\Ii$, and to define
$\Ss=\gr_\Jj(\Oo_{\Xx'})$, and consider $\gr_\Ii(\formelF)$. In
any event, \Ref{IX.2.1} can be applied directly to the
morphism $\widehat{f}:\widehat{X}\to S$ considered in the present
subsection, where one simply takes $\Jj=0$. Thus, to verify that
$\widehat{f}_\ast(\formelF)$ is coherent, it suffices, by virtue of
\loccit, to verify that $\R^if_{0\ast}(\gr_\Ii(\formelF))$ is
coherent on $S$ for $i=0, 1$; for this, note that by virtue of
a) and b), the Module in question is none other than
$\sisi{\coprod}{\bigoplus}_{m\geq 0}\R^if_{0\ast}(F_0(-m))$, which
is indeed coherent by virtue of hypothesis c) and
\Ref{XII.1.5}.

This proves \Ref{XII.3.2}~(i). For \Ref{XII.3.2}~(ii), we shall need the

\begin{lemme} \label{XII.3.3}
Under\pageoriginaled the conditions of \Ref{XII.3.1a}), \Ref{XII.3.1b}), \Ref{XII.3.1c}) of \Ref{XII.3.1}, set
\[
G_m= \widehat{f}_\ast(\formelF(\cdot)_m) = \sisi{\coprod}{\tbigoplus}_n\widehat{f}_\ast(\formelF_m(n))
\]
Then the projective system $(G_m)$ satisfies the Mittag-Leffler condition.
\end{lemme}

We may suppose that $S$ is affine, with ring $A$. Let $\Ss$ be a finitely generated graded $A$-algebra in positive degrees, and let $t'\in\Ss_1$, such that $X$ embeds in $\Proj(\Ss)$, with $\OX(1)$ induced by \sisi{$\SheafProj(\Ss(1))$}{$\Oo(1)$} and the section $t$ the image of $t'$. Endow $\Ss$ with the $\Jj$-adic filtration, where $\Jj=t'\Ss$, and consider the projective system of the $\formelF(\cdot)_m$ in the category of abelian sheaves on $X_0$. We are again under the preliminary conditions of $\EGA 0_{\textup{III}}$ 13.7.7\sfootnote{Corrected as indicated in \Ref{IX} p\ptbl\sisi{11}{\pageref{pagerectif}.}}, and moreover $\H^i(X_0, \gr(\formelF(\cdot))$ is a finite-type Module over $\gr_\Jj(\Ss)$ for $i=0, 1$. Indeed, since $t'$ is $\formelF$-regular, one sees at once that, as a Module over $(\Ss /t'\Ss)[T]$ (of which $\gr_\Jj(\Ss)$ is a quotient), the Module in question identifies with $\H^i(X_0, F_0(\cdot))\otimes_{\Ss /t'\Ss}(\Ss /t'\Ss)[T]$; but by \Ref{XII.1.5}, $\H^i(X_0, F_0(\cdot))$ is of finite type over $\Ss$, and hence over $\Ss /t'\Ss$, for $i=0, 1$, which proves our assertion. Consequently, we are under the conditions for applying $0_{\text{III}}$ 13.7.7 with $n=1$, which proves \Ref{XII.3.3}.

This point being established (and still supposing $S$ affine, as we may for proving \Ref{XII.3.2}~(ii)), let $m_0$ be such that $m\geq m_0$ implies $\Im(G_m\to G_0)= \Im(G_{m_0}\to G_0)$, so that both sides are also equal to $\Im(\varprojlim G_m\to G_0) = \Im(\widehat{f}_\ast(\formelF(\cdot))\to f_\ast\formelF_0(\cdot))$. Now note that for large $n$, $\formelF_{m_0}(n)$ is generated by its sections, and therefore $\formelF_0(n)$ is generated by sections that lift to $\formelF_{m_0}$, and hence (by the choice of~$m_0$) that lift to $\formelF$. Thus the sections of $\formelF$ generate $\formelF_0$, and hence also~$\formelF$ by Nakayama. This proves \Ref{XII.3.2}~(ii), and therefore \Ref{XII.3.1}.

\begin{corollaire} \label{XII.3.4}
Let $f:X\to S$ be a projective flat morphism, with $S$ locally noetherian, let $\OX(1)$ be an invertible Module on $X$, relatively ample over $S$, and let $t$ be a section of this Module such that for every $s\in S$, the induced section $t_s$ on\pageoriginale the fiber $X_s$ is $\Oo_{X_s}$-regular; let $X_0$ be the closed subprescheme of zeros of $t$, and let $\widehat{X}$ be the formal completion of $X$ along~$X_0$. Suppose that for every $s\in S$, $X_{0s}$ has depth $\geq 2$ at its closed points (\ie $X_s$ has depth $\geq 3$ at the closed points of $X_{0s}$), and $X_s$ has depth $\geq 2$ at its closed points. Under these conditions, the pair $(X, X_0)$ satisfies the effective Lefschetz condition ($\Leff$) \sisi{of \Ref{X}~\Ref{X.2}}{of paragraph~\Ref{X.2} of Expos\'e~\emph{\Ref{X}}}, \ie:
\begin{enumeratea}
\item
For every open neighborhood $U$ of $X_0$ in $X$, the functor
\[
F \sisi{\rightsquigarrow}{\mto} \widehat{F}
\]
from the category of coherent locally free Modules on $U$ to the category of coherent locally free Modules on $\widehat{X}$ is fully faithful.
\item
For every coherent locally free Module $\formelF$ on $\widehat{X}$, there exist an open neighborhood~$U$ of $X_0$ and a coherent locally free Module $F$ on $U$ such that $\formelF$ is isomorphic to~$\widehat{F}$.
\end{enumeratea}
\end{corollaire}

Indeed, a) has already been noted in \Ref{XII.2.4} under weaker conditions. For b), we apply \Ref{XII.3.1}, which gives the conclusion, at least if $S$ is noetherian and admits an absolutely ample invertible Module, in particular if $S$ is affine. Indeed, if $F$ is a coherent module on $X$ such that $\widehat{F}$ is isomorphic to $\formelF$, and hence locally free, it follows that $F$ is locally free on a neighborhood $U$ of $X_0$, and $\sisi{F|U}{F|_U}$ will satisfy the required condition. But now note that by \Ref{XII.2.5}, for such an $F$, its direct image under the immersion $U\to X$ is coherent, and is moreover independent of the chosen solution $(U, F)$ (taking account \sisi{\ignorespaces}{of the fact} that two solutions coincide in a neighborhood of $X_0$, by a)). More precisely, one can find a coherent Module $F$ on $X$ and an isomorphism $\widehat{F}\isomto \formelF$, such that $F$ has depth $\geq 2$ at all the points of $X$ that are closed in their fiber, and this determines $F$ up to unique isomorphism. By this uniqueness property, the solutions of the problem obtained over the affine open subsets of $S$ glue together, giving a coherent $F$ on all of $X$ and an isomorphism $\widehat{F}\simeq \formelF$. Restricting $F$ to the open subset $U$ of points at which it is free gives what was required.

Thanks\pageoriginale to \Ref{XII.2.4} and \Ref{XII.3.4}, in the setting of a projective algebraic scheme and a \og hyperplane section\fg thereof, one can apply the general facts established in Expos\'es \Ref{X} and \Ref{XI} concerning conditions (Lef) and (Leff). Thus:

\begin{corollaire} \label{XII.3.5} Let $X$ be a projective algebraic scheme equipped with an ample invertible module $\OX(1)$, let $t$ be a section of this module that is $\OX$-regular, and let $X_0$ be the closed subscheme of zeros of $t$. Assume that $X$ has depth $\geq 2$ at its closed points (\resp and depth $\geq 3$ at the closed points of $X_0$). Then $\pi_0(X_0)\to \pi_0(X)$ is bijective; in particular, $X$ is connected if and only if $X_0$ is connected, and, after choosing a geometric base point in $X_0$, $\pi_1(X_0)\to \pi_1(X)$ is surjective. More generally, for every open $U\supset X_0$, the homomorphism $\pi_1(X_0)\to \pi_1(U)$ is surjective (\resp the homomorphism $\pi_1(X_0)\to \varprojlim_U\pi_1(U)$ is bijective). In the parenthesized case, if one assumes in addition that the local ring at every closed point of $X$ not in $X_0$ is pure (\Ref{X.3.2}) (for example, regular, or merely a complete intersection), then $\pi_1(X_0)\to \pi_1(X)$ is an isomorphism.
\end{corollaire}

Apply \Ref{X.2.4} and \Ref{XII.3.4}. Note that, in the parenthesized case, the hypothesis that $X$ have depth $\geq 3$ at the closed points of $X_0$ implies that every irreducible component of $X$ having dimension $\neq 0$ has dimension $\geq 3$ (as one sees by noting that such a component necessarily meets $X_0$, and by considering a closed point of the intersection).

\begin{remarquestar}
When $X$ is normal and has dimension $\geq 2$ at all its points, it has depth $\geq 2$ at its closed points, and the non-parenthesized hypotheses of \Ref{XII.3.5} hold. In this case there is a more elementary proof of the surjectivity of $\pi_1(X_0)\to \pi_1(X)$ using Bertini's theorem (\cf \SGA 1 X.2.10). If one assumes in addition that $X_0$ is normal and that $X$ has dimension $\geq 3$ at all its points, then the parenthesized hypotheses of \Ref{XII.3.5} hold. In this case, \Ref{XII.3.5} was established by Grauert (indeed, thanks to the normality hypothesis, one can then dispense with the existence theorem \Ref{XII.3.1} by various devices)\pageoriginale. This proof of Grauert's was the starting point of the \og Lefschetz theory\fg that is the subject of the present seminar.
\end{remarquestar}

\begin{corollaire} \label{XII.3.6} Let $X$, $\OX(1)$, $t$, and $X_0$ be as in \Ref{XII.3.5}. Assume that $X$ has depth $\geq 2$ at its closed points, and that
\[
\H^i(X_0, \Oo_{X_0}(-n))=0
\]
for $n>0$ and for $i=1$ (\resp for $i=1$ and $i=2$), which, by \Ref{XII.1.4}, implies that $X_0$ has depth $\geq 2$ (\resp $\geq 3$) at its closed points, \ie that $X$ has depth $\geq 3$ (\resp $\geq 4$) at the closed points of $X_0$. Under these conditions, for every open neighborhood $U$ of $X_0$, $\Pic(U)\to\Pic(X_0)$ is injective; in particular, $\Pic(X)\to\Pic(X_0)$ is injective (\resp $\varprojlim_U \Pic(U)\to\Pic(X_0)$ is bijective). In the parenthesized case, if one assumes in addition that the local ring of $X$ at every closed point not in $X_0$ is parafactorial (\Ref{XI.3.1}) (for example, regular, or more generally a complete intersection), then $\Pic(X)\to\Pic(X_0)$ is bijective.
\end{corollaire}
Apply \Ref{XI}~\Ref{XI.3.12} and \Ref{XI.3.13}, noting that the parenthesized hypothesis implies that the irreducible components of $X$ having dimension $\neq 0$ have dimension $\geq 4$. In particular, applying this when $X$ is a global complete intersection of dimension $\geq 4$ in projective space, one obtains:

\begin{corollaire} \label{XII.3.7} Let $X$ be an algebraic scheme of dimension $\geq 3$ that is a complete intersection in the scheme $\PP_k^r$. Then $\Pic(X)$ is the free group generated by the class of the sheaf $\OX(1)$.
\end{corollaire}
Proceed by induction on the number of hypersurfaces whose intersection is $X$, applying \Ref{XII.3.6} and noting that, for a complete intersection $X$ of dimension $\geq 3$, one has $\H^i(X, \OX(n))=0$ for $i=1,2$ and every $n$.

\begin{remarque} \label{XII.3.8}
In\pageoriginale the case where $X$ is a nonsingular hypersurface, \Ref{XII.3.7} is due to Andreotti. The result \Ref{XII.3.7} can also be expressed (when $X$ is nonsingular) by saying that the homogeneous coordinate ring of $X$ is factorial, and in this form it is contained in \Ref{XI}~\Ref{XI.3.13} (ii). We also note that Serre had given a proof of \Ref{XII.3.7} in the nonsingular case by transcendental methods, using a specialization argument to reduce to characteristic 0, where the Lefschetz theorem is available in its classical form. Of course, the fact that the purely algebraic proof given here makes it possible to dispense with nonsingularity hypotheses in the statement of the Lefschetz theorem invites one to reconsider it in the classical case as well. \Cf the following expos\'e, which proposes conjectures in this direction.

In Corollaries \Ref{XII.3.5} and \Ref{XII.3.7}, we have worked over a base field, whereas the key theorems \Ref{XII.2.4} and \Ref{XII.3.4} are valid over an arbitrary base. To generalize Corollaries \Ref{XII.3.5} and \Ref{XII.3.6} to a general $S$, we must give workable criteria for a point of $X$ (flat over $S$) to have a \og pure\fg \resp parafactorial local ring. This will be the subject of the following n\textsuperscript{o}.
\end{remarque}

\section{Formal completion and normal flatness} \label{XII.4}

\setcounter{equation}{18}

\begin{theoreme} \label{XII.4.1} Let $X$ be a locally Noetherian prescheme, locally embeddable in a regular scheme, $Y$ a closed subset of $X$, $U=X- Y$, $X_0$ a closed subprescheme of $X$ defined by an ideal $\Jj$, $\widehat{X}$ the formal completion of $X$ along $X_0$, $U_0$ the trace of $X_0$ on $U$, $\widehat{U}$ the formal completion of $U$ along $U_0$, $i:U\to X$ and $\widehat{i}:\widehat{U}\to \widehat{X}$ the canonical immersions, and $n$ an integer. Assume that
\begin{enumeratea}
\item
$X$ is normally flat along $X_0$ at the points of $Y\cap X_0$, \ie at these points the \sisi{m}{M}odules $\Jj^n/\Jj^{n+1}$ over $X_0$ are flat, \ie \sisi{\ignorespaces}{locally} free.
\item
For every $x\in Y\cap X_0$, one has $\prof \Oo_{X_0, x}\geq n+2$.
\end{enumeratea}
Under\pageoriginale these conditions, the following assertions hold:
\begin{enumerate}
\sisi{\item[$1^\circ)$]}{\item} Let $F$ be a coherent Module on $U$, and suppose that:
\begin{enumerate}
\item[\textup{c)}] For every $x\in Y- Y\cap X_0$, one has $\prof \Oo_{X, x}\geq n+2$.

\item[\textup{d)}] $F$ is free at the points of $U_0$, and has depth $\geq n+1$ at all the points of $U$ where it is not free.
\end{enumerate}
Then the graded Module
\[
\sisi{\coprod}{\tbigoplus}_{m\geq 0}\R^pi_\ast(\Jj^mF)
\]
over $\sisi{\coprod}{\bigoplus}_{m\geq 0}\Jj^m$ is of finite type for $p\leq n$. \sisi{\item[$2^\circ$]}{\item} Let $\formelF$ be a coherent Module on $\widehat{U}$; then the graded Module
\[
\sisi{\coprod}{\tbigoplus}_{m\geq 0}\R^pi_\ast(\Jj^m\formelF/\Jj^{m+1}\formelF)
\]
over $\sisi{\coprod}{\bigoplus}_{m\geq 0}\Jj^m/\Jj^{m+1}=\gr_{\Jj}(\OX)$ is of finite type for $p\leq n$.
\end{enumerate}
\end{theoreme}

\begin{proof}
1) Let $X'\!=\!\Spec(\sisi{\coprod}{\bigoplus}_{m\geq 0}\Jj^m)$; the base change \hbox{$f:X'\!\to\! X$} then defines $U'=X'- Y'$, $X_0'$, $U_0'=X_0'\sisi{\cap}{} U'$, and immersions $i':U'\to X'$, \hbox{$i_0':U_0'\to X_0'$}. Thus we have a Cartesian square
\[
\xymatrix{
X'\ar[d]_f&\ar[l]_-{i'} U'\ar[d]^g \\
X&\ar[l]_-{i} U
}
\]
and we have
\steco
\begin{equation} \label{eq:XII.19}
\sisi{\coprod}{\tbigoplus}_{m\geq 0}\R^pi_\ast(\Jj^mF) = \R^pi_\ast \Big(\sisi{\coprod}{\tbigoplus}_{m\geq 0} \Jj^mF \Big) =\R^pi_\ast (g_\ast(F')),
\end{equation}
where\pageoriginale $F'=g^\ast F$, so that there is indeed a canonical isomorphism
\[
g_\ast(F') \isomto \sisi{\coprod}{\tbigoplus}_{m\geq 0}\Jj^mF,
\]
since this is true at the points of $U_0$, because $F$ is free there by d), and also at the points of $U_0$, because there one has $\Jj^m=\Oo_U$, (so that in both cases, $\Jj^m\otimes_{\Oo_U}F\to \Jj^mF$ is an isomorphism).

On the other hand, since $f$, and hence $g$, are affine, one has
\steco
\begin{equation} \label{eq:XII.20} {}\R^pi_\ast(g_\ast(F')) = \R^p(ig)_\ast(F') = \R^p(fi')_\ast(F') = f_\ast(R^pi'_\ast(F'))
\end{equation}
and therefore, comparing (\Ref{eq:XII.19}) and (\Ref{eq:XII.20}), we see that \hyperref[XII.4.1]{assertion 1}\sisi{${}^\circ$}{} is equivalent to the following: $R^pi'_\ast(F')$ is a Module of finite type, \ie coherent, on $X'$, for every $p\leq n$. Now, since $X$ is locally embeddable in a regular scheme, the same is true of $X'$, which is of finite type over $X$, and we may apply the coherence criterion \sisi{\ignorespaces}{\Ref{VIII}}~\Ref{VIII.2.3} to a coherent extension $F^{''}$ of $F'$: we wish to express that $\SheafH^p_Y(F^{''})$ is coherent for $p\leq n+1$, and this is also equivalent to saying that, for every $x'\in U'$ such that
\begin{equation*} \label{eq:XII.20bis} \tag{20 bis} {}\codim (\overline{x'}\cap Y', \overline{x'})=1,
\end{equation*}
one has
\steco
\begin{equation} \label{eq:XII.21} {} \prof F'_{x'}\geq n+1.
\end{equation}
Now this condition is satisfied at the points $x'$ where $F'$ is not free, because for such an $x'$ one has $x'\notin U_0'$ by d), and therefore $g$ is an isomorphism there; again by d), $F$ has depth $\geq n+1$ at $g(x')$, and hence $F'$ has depth $\geq n+1$ at $x'$. It is therefore enough to verify condition (\Ref{eq:XII.21}) at those $x'\in U'$ satisfying (\Ref{eq:XII.21}) at which $F'$ is free. For this, it is enough to prove that one has\pageoriginale
\begin{equation*} \label{eq:XII.21bis} \tag{21 bis}
\prof \Oo_{X', x'}\geq n+1
\end{equation*}
\enlargethispage{\baselineskip}%
at these points; \afortiori it is enough to establish this relation at all the points~$x$ of $U'$ satisfying (\Ref{eq:XII.20}\sisi{}{ bis}). Now, \sisi{by the criterion cited}{again by criterion~\Ref{VIII.2.3}} in Expos\'e \Ref{VIII}, this is equivalent to the assertion that the Modules
\[
\SheafH^p_{Y'}(\Oo_{X'}) \hspace{5mm} \text{for} \; p\leq n+1
\]
are coherent. In fact, we shall prove that they are even zero, or, equivalently by Expos\'e~\Ref{III}, that one has
\steco
\begin{equation} \label{eq:XII.22} {} \prof \Oo_{X', x'}\geq n+2 \hspace{5mm} \text{for every }\; x'\in Y'.
\end{equation}
For this, we distinguish two cases. If $x'\notin X'_0$, then $f$ is an immersion at $x'$, and we must verify that $F$ has depth \sisi{at}{$\geq n+2$ at } the image $x=f(x')$, which is precisely condition c). If, on the other hand, $x'\in X'_0$, \ie $x=f(x')\in X_0$ and hence $x\in Y\cap X_0$, we apply conditions a) and b) by means of the

\begin{lemme} \label{XII.4.2}
Let $X$ be a locally Noetherian prescheme, $X_0$ a closed subprescheme of $X$ defined by an ideal $\Jj$, $X'=\Spec(\sisi{\coprod}{\bigoplus}_{m\geq 0} \Jj^m)$, $X_0'=\Spec(\sisi{\coprod}{\bigoplus}_{m\geq 0} \Jj^m/\Jj^{m+1}) = X'\times_XX_0$, and $x$ a point of $X_0$ at which $X$ is normally flat along $X_0$, \ie such that $\gr_{\Jj}(\OX)= \sisi{\coprod}{\bigoplus}_{m\geq 0}\Jj^m/\Jj^{m+1}$ is flat there as a Module over $X_0$. Then, for every sequence of elements $f_i$ ($1\leq i\leq m$) of $\Oo_{X, x}$ whose images in $\Oo_{X_0, x}$ form an $\Oo_{X_0, x}$-regular sequence, and every $x'\in X'$ above $x$, the images of the $f_i$ in $\Oo_{X', x'}$ (respectively in $\Oo_{X'_0, x'}$) likewise form an $\Oo_{X', x'}$-regular sequence (respectively an $\Oo_{X'_0, x'}$-regular sequence); in particular, one has
\steco
\begin{equation} \label{eq:XII.23}
\prof \Oo_{X', x'} \geq \prof \Oo_{X_0, x}, \quad \prof \Oo_{X'_0, x'} \geq \prof \Oo_{X_0, x}.
\end{equation}
\end{lemme}

To\pageoriginale prove it, we may suppose that $X$ is local with closed point $x$, and hence affine with ring $A=\Oo_{X, x}$, with $\Jj$ defined by the Ideal $J$; it is enough to prove that, for every sequence $f_i$ ($1\leq i\leq m$) of elements of $A$ whose images in $A/J$ form an $A/J$-regular sequence, the $f_i$ likewise form a $(\sisi{\coprod}{\bigoplus}_{m\geq 0} J^m)$-regular sequence and a $(\sisi{\coprod}{\bigoplus}_{m\geq 0} J^m/J^{m+1})$-regular sequence, \ie for every $m$, it forms a $J^m$-regular sequence and a $J^m/J^{m+1}$-regular sequence. The second assertion is trivial, since $J^m/J^{m+1}$ is a free module over $A/J$. The first follows by considering the $J$-adic filtration of $J^m$ and noting that the sequence of the~$f_i$ is regular on the graded module associated with $J^m$ for that filtration.

This proves \Ref{XII.4.2}, and consequently \Ref{XII.4.1}~1\sisi{\noo}{}).

Let us prove \Ref{XII.4.1} 2\sisi{\noo}{}). For this, use the Cartesian square
\[
\xymatrix{
X'_0\ar[d]_{f_0}&\ar[l]_-{i'_0} U'_0\ar[d]^{g_0} \\
X_0&\ar[l]_-{i_0} U_0
}
\]
and, proceeding as at the beginning of the proof of 1\sisi{\noo}{}), we find that
\steco
\begin{equation} \label{eq:XII.24}
\sisi{\coprod}{\tbigoplus}_{m\geq 0} \R^pi_\ast(\Jj^m\formelF/\Jj^{m+1}\formelF) \simeq {f_0}_\ast (\R^p{\sisi{i}{i'}_0}_\ast(\formelF_0')),
\end{equation}
where $\formelF_0=\formelF/\Jj\formelF$ and $\formelF_0'=g_0^\ast(\formelF_0)$, (using the fact that $\formelF$ is locally free). Thus the conclusion of 2\sisi{\noo}{}) is equivalent to saying that, for $p\leq n$, $\R^p{i'_0}_\ast(\formelF_0')$ is a coherent Module.

\enlargethispage{\baselineskip}%
Here again, taking into account that $\formelF_0'$ is locally free, criterion \Ref{VIII}~\Ref{VIII.2.3} allows us to reduce to proving the same assertion after replacing $\formelF_0'$ by $\Oo_{U'_0}$, \ie to proving that the Modules\pageoriginale
\[
\SheafH^p_{Y_0'}(\Oo_{X_0'}) \hspace{5mm} \text{for} \; p\leq n+1 \hspace{5mm} \text{(where} \; Y'_0=Y'\cap X_0'=X_0'- U'_0 \text{)}
\]
are coherent. Once again, we prove that they are in fact zero, i.e., that one has
\steco
\begin{equation} \label{eq:XII.25} {}\prof \Oo_{X'_0, x'} \geq n+2 \hspace{5mm} \text{for every} \; x'\in Y'_0.
\end{equation}
This does indeed follow from conditions a) and b), in view of \Ref{XII.4.2}. This completes the proof of \Ref{XII.4.1}.
\skipqed
\end{proof}

\begin{remarque} \label{XII.4.3}
It follows immediately, by descent, that the hypothesis that $X$ be locally embeddable in a regular scheme can be replaced by the following weaker hypothesis: there exists a faithfully flat and quasi-compact morphism $\overline{X}\to X$ such that $\overline{X}$ is locally embeddable in a regular scheme.
\end{remarque}

Theorem \Ref{XII.4.1} enables us to apply the results of Expos\'e \Ref{IX} (the comparison and existence theorems). We shall be interested in particular in the

\begin{corollaire} \label{XII.4.4}
Assume that conditions a), b), and c) of Theorem \Ref{XII.4.1} hold, with $n=1$ and $X=\Spec(A)$, where $A$ is separated and complete for the $J$-adic topology. Then
\begin{enumerate}
\sisi{\item[\textup{1\noo})]}{\item} The functor $F\mto \widehat{F}$ from the category of locally free coherent modules on~$U$ to the category of locally free coherent modules on $\widehat{U}$ is fully faithful. \sisi{\item[\textup{2\noo})]}{\item} For every locally free coherent module $\formelF$ on $\widehat{U}$, there exists a coherent module $F$ on $U$ and an isomorphism $\widehat{F}\isomto \formelF$.
\end{enumerate}
In particular, if, for every $x\in U$ whose closure in $U$ does not meet $U_0$, \ie such that $\overline{x}\cap X_0\subset Y$, one has $\prof \Oo_{U, x}\geq 2$, then the pair $(U, U_0)$ satisfies the effective Lefschetz condition ($\Leff$) of Expos\'e \Ref{X}.
\end{corollaire}
\noindent
(For the last assertion, one proceeds as in \Ref{X}~\Ref{X.2.1}.)

\pagebreak[2]
Special case of \Ref{XII.4.4}:

\begin{corollaire} \label{XII.4.5}
Let\pageoriginaled $A$ be a Noetherian ring, let $J$ be an ideal of $A$ contained in the radical, and let $A_0=A/J$. Assume
\begin{enumeratei}
\item
$\prof A_0\geq 3$.
\item
$\gr_J(A)$ is a free $A_0$-module.
\item
$A$ is complete for the $J$-adic topology.
\end{enumeratei}

Let $X=\Spec(A)$, $X_0=\Spec(A_0)=V(J)$, let $a$ be the closed point of $X$, let $U=X-\{a\}$, $U_0=X_0-\{a\}$, and let $\widehat{U}$ be the formal completion of $U$ along $U_0$. Then the functor $F\mto \widehat{F}$ from the category of locally free coherent modules on $U$ to the category of locally free coherent modules on $\widehat{U}$ is fully faithful. Moreover, for every locally free coherent module $\formelF$ on $\widehat{U}$, there exists a coherent module (not necessarily locally free!) $F$ on $U$, and an isomorphism $\widehat{F}\simeq\formelF$.
\end{corollaire}

Note that, thanks to \Ref{XII.4.3}, we did not have to assume that $A$ is a quotient of a regular ring, because the completion of $A$ for the $\rr(A)$-adic topology satisfies this condition in any case.

Proceeding as in Expos\'es \Ref{X} and \Ref{XI}, we deduce from \Ref{XII.4.5}

\begin{corollaire} \label{XII.4.6}
Under the conditions of \Ref{XII.4.5}, one has the following:
\begin{enumeratea}
\item
$U$ and $U_0$ are connected (\Ref{III}~\Ref{III.3.1}).

Choosing a geometric base point in $U_0$, the homomorphism
\[
\pi_1(U_0)\to \pi_1(U)
\]
is surjective.
\item
The homomorphism
\[
\Pic(U)\to \Pic(U_0)
\]
is injective.
\end{enumeratea}
\end{corollaire}

To\pageoriginale prove b), in view of \Ref{XII.4.5}, this amounts to checking that every isomorphism $L_0'\isomto L_0$ lifts to an isomorphism $\widehat{L'}\isomto \widehat{L}$. For this, one lifts step by step to isomorphisms $L_n'\isomto L_n$; the obstructions lie in $\H^1(U_0, \Jj^n/\Jj^{n+1})$, and these modules are zero because $J^n/J^{n+1}$ is free and $\prof A_0\geq 3$.

We are now in a position to prove the
\enlargethispage{-1\baselineskip}%
\begin{theoreme} \label{XII.4.7}
Let $A$ be a noetherian local ring, let $J$ be an ideal of $A$ contained in its radical, and let $A_0=A/J$. Suppose that
\begin{enumeratei}
\item
$\prof A_0\geq 3$.
\item
$\gr_J(A)$ is a free module over $A_0$.
\end{enumeratei}
\noindent
Then, if $A_0$ is \og pure\fg (\Ref{X}~\Ref{X.3.1}) (\resp parafactorial (\Ref{XI}~\Ref{XI.3.1})), the same holds for~$A$.
\end{theoreme}

\begin{proof}
By descent, we may suppose that we also have
\begin{enumeratei}\setcounter{enumi}{2}
\item
$A$ is complete for the $J$-adic topology.
\end{enumeratei}

Indeed, by (i) and (ii), we have $\prof(A)\geq 3$, and hence $\prof(\widehat{A})\geq 3$, where $\widehat{A}$ is the completion of $A$ for the $J$-adic topology, and we apply \Ref{X}~\Ref{X.3.6} and \Ref{XI}~\Ref{XI.3.6}. We are therefore under the conditions of \Ref{XII.4.5}. Since $\prof(A)\geq 3\geq 2$, to say that $A$ is parafactorial means simply that $\Pic(U)=0$, and by \Ref{XII.4.6}~b), it is enough for this that $\Pic(U_0)=0$, \ie that $A_0$ be parafactorial. To prove that $A$ is \og pure\fg if $A_0$ is, we must prove that if $V$ is an \'etale covering of $U$, defined by an algebra $\Bb$ on~$U$, then $\H^0(U, \Bb)$ is a finite \'etale algebra over $A$. Now, since $A_0$ is pure, the same holds for the $A_n$ (which differ from it only by nilpotent elements), and therefore, for every $n$, $\sisi{}{B_n=}\H^0(U, \Bb_n)$ is an \'etale algebra over $A/J^{n+1}$; these algebras of course fit together, so that $\varprojlim B_n$ is an \'etale algebra over $A$. By \Ref{XII.4.5}, this algebra is none other than $\H^0(U, \Bb)$, which proves our assertion.
\skipqed
\end{proof}

\begin{corollaire} \label{XII.4.8}
Let\pageoriginale $f:X\to Y$ be a flat morphism of locally noetherian preschemes, let $x\in X$, and let $y=f(x)$. Suppose that $\Oo_{X_y, x}$ is a \og pure\fg (\resp parafactorial) local ring of depth $\geq 3$. Then the same holds for $\Oo_{X, x}$.
\end{corollaire}

This is the result of the kind promised at the end of the preceding \sisi{N\textsuperscript{o}}{\numero}, for generalizing Corollaries \Ref{XII.3.5} and following. Thus, using \Ref{XII.3.4}, we obtain the

\begin{corollaire} \label{XII.4.9}
Let $f:X\to S$ be a projective flat morphism, with $S$ locally Noetherian, let $\OX(1)$ be an invertible module on $X$, ample relative to $S$, let $t$ be a section of $\OX(1)$ such that, for every $s\in S$, the section $t_s$ induced on $X_s$ is $\Oo_{X_s}$-regular, let $X_0$ be the closed subscheme of zeros of $t$, and let $X_m$ be the closed sub-\sisi{pre}{}scheme of zeros of $t^{m+1}$. Assume that, for every $s\in S$, $X_s$ has depth $\geq 3$ at all its closed points. Then:
\begin{enumeratea}
\item
If the local rings at the closed points of $X_s-X_{0,s}$ ($s\in S$) are \og pure\fg\,---for example, if they are complete intersections---then the functor $X'\mto X_0'=X'\times_X X_0$ from the category of finite \`etale coverings of $X$ to the category of finite \`etale coverings of $X_0$ is an equivalence of categories; in particular, after choosing a geometric base point in $X_0$, the homomorphism
\[
\pi_1(X_0)\to \pi_1(X)
\]
is an isomorphism\refstepcounter{toto}\nde{\label{FH}Let us point out
the striking connectedness result obtained subsequently by Fulton and
Hansen, in the case where $S=\Spec(k)$ ($k$ an algebraically closed
field). Let $g:X\to {\PP^m_k}\times{\PP^m_k}$ be such that $\dim
g(X)>m$; then the inverse image of the diagonal is connected. Among
other things, this makes it possible to generalize
Corollary~\Ref{XII.4.9} when $f$ is the structural morphism
of projective space $\PP^m_k$ over $S=\Spec(k)$: more precisely, an
irreducible subvariety $X$ of ${\PP^m_k}$ of dimension $>m/2$ has
trivial fundamental group! (\cf Fulton~W. \& Hansen J., {\og
A connectedness theorem for projective varieties, with
applications to intersections and singularities of mappings\fg},
\emph{Ann. of Math. (2)} \textbf{110} (1979), \numero 1, p\ptbl
159-166). For generalizations to Grassmannians or
abelian varieties, see {Debarre~O.}, {\og Th\'eor\`emes de
connexit\'e pour les produits d'espaces projectifs et les
grassmanniennes\fg}, \emph{Amer. J.~Math.} \textbf{118} (1996),
\numero 6, p\ptbl 1347-1367 and {\og Th\'eor\`emes de connexit\'e et
vari\'et\'es ab\'eliennes\fg}, \emph{Amer. J.~Math.} \textbf{117}
(1995), \numero 3, p\ptbl 787--805. The triviality result for the
fundamental group of $X$ above was obtained independently
by Faltings, who moreover proves that
$\text{Pic}(X)$ has no torsion prime to the characteristic
of $k$, using methods of algebraization of formal vector bundles,
more in the spirit of Grothendieck's techniques, \cf
(Faltings~G., \og{Algebraization of some formal vector bundles\fg}
\emph{Ann. of Math. (2)} \textbf{110} (1979), \numero 3, p\ptbl
501--514).}

\item
If the local rings at the closed points of $X_s-X_{0,s}$ ($s\in S$) are \og parafactorial\fg---for example, regular, or complete intersections of dimension $\geq 4$---then, for every integer $m$ such that $\R^if_{0\ast}(\Oo_{X_0}(-n))=0$ for $n>m$ and $i=1,2$, the map $\Pic(X)\to \Pic(X_m)$ is bijective.

Moreover, if $S$ is Noetherian and the $X_{0,s}$ have depth $\geq 3$ at their closed points, such $m$ exist (\cf \Ref{XII.1.5}).
\end{enumeratea}
\end{corollaire}

\begin{remarque} \label{XII.4.10}
Under\pageoriginale the conditions of the last assertion of \Ref{XII.4.9} b), we saw in \Ref{XII.1.5} that there exists an $m$ such that $n>m$ even implies $\H^i(X_0, \Oo_{X_0}(-n))=0$ for $i=1,2$, and indeed for $i\leq 2$). This condition is stronger than $\R^if_{\ast}(\Oo_{X_0}(-n))=0$ for $i=1,2$, and it has the further advantage of being stable under base change. The same is true of the depth hypotheses made in \Ref{XII.4.9}, and also of a hypothesis of the type \og the $X_s$ are locally complete intersections\fg. It follows under these conditions that \Ref{XII.4.9} b) also implies that the morphism of functors
\[
\SheafPic_{X/S}\to \SheafPic_{X_n/S}
\]
in $\catSch_{/S}$ is an isomorphism, and hence so is the morphism of relative Picard schemes, when these exist:
\[
\Pic_{X/S}\to \Pic_{X_n/S}.
\]
Even when $S$ is the spectrum of an algebraically closed field, this statement is considerably more precise than the statement that merely says that $\Pic(X)\to \Pic(X_n)$ is bijective.

One may ask whether one can always take $n=0$ in the preceding conclusions (thus assuming that the $X_{0,s}$ have depth $\geq 3$ at their closed points). When $X_0$ is smooth over $S$ and the residual characteristics of $S$ are zero, this is indeed so, by the Kodaira \og vanishing theorem\fg (proved by transcendental methods, using a K\"ahler metric), which implies that, for every connected smooth projective scheme of dimension $n$ over a field $k$ of characteristic zero and every ample invertible module $L$ on $X$, one has $\H^i(X,L^{-1})=0$ for $i\neq n$. It is not known\nde{As Raynaud observed, the decomposition result for the de Rham complex of Deligne and Illusie readily implies the vanishing of the group $H^i(X,L^{-1})$ (with $L$ ample on $X$, projective and smooth over $k$ of characteristic $p>0$) for $i<\inf(p,\dim(X))$, provided that $X$ can be lifted to a flat scheme over $W_2(k)$ (\cf Deligne~P. \& Illusie L., {\og Rel\`evements modulo~$p^2$ et d\'ecomposition du complexe de de Rham\fg}, \emph{Invent. Math.} \textbf{89} (1987), \numero 2, p\ptbl 247--270); this gives a purely algebraic proof of Kodaira's result for projective varieties in characteristic zero. If $X$ cannot be lifted, it is well known that the \og vanishing theorem\fg is false; \cf the example in (Raynaud~M., {\og Contre-exemple au "vanishing theorem" en caract\'eristique $p>0$\fg}, in \emph{C.P.~Ramanujam---a tribute}, Tata Inst. Fund. Res. Studies in Math., vol.~8, Springer, Berlin-New York, 1978, p\ptbl 273--278); see also the very elegant examples in (Haboush~W. \& Lauritzen~N., {\og Varieties of unseparated flags\fg}, in \emph{Linear algebraic groups and their representations (Los Angeles, CA, 1992)}, Contemp. Math., vol.~153, American Mathematical Society, Providence, RI, 1993, p\ptbl 35--57), simplified in (Lauritzen~N. \& Rao~A.P., {\og Elementary counterexamples to Kodaira vanishing in prime characteristic\fg}, \emph{Proc. Indian Acad. Sci. Math. Sci.} \textbf{107} (1997), \numero 1, p\ptbl 21--25). On the other hand, I know of no example in which the map $\Pic(X_{n+1})\to\Pic(X_n)$ is not surjective for $n>1$ in positive characteristic, where $X_n$ denotes a thickened hyperplane section of a smooth projective $X$ as above.} at present whether this theorem can be replaced by a generalization in characteristic $p>0$, and whether the smoothness hypothesis can be replaced by a hypothesis of a more general nature (involving depth, or of \og complete intersection\fg type, and so on).
\end{remarque}

\section[Universal finiteness conditions]
{Universal finiteness conditions for a nonproper morphism} \label{XII.5} Recall\pageoriginale for reference the \setcounter{equation}{25}

\begin{proposition} \label{XII.5.1}
Let $f:X\to S$ be a proper morphism of preschemes, with $S$ locally noetherian, let $U$ be an open subset of $X$, let $g:U\to X$ be the canonical immersion, let $h=fg:U\to S$, and let $F$ be a Module on $U$. Suppose that the Modules $\R^ig_\ast(F)$ are coherent for $i\leq n$ (a hypothesis local in nature on $X$, which in practice is checked using the criterion \Ref{VIII}~\Ref{VII.2.3}). Then $\R^ih_\ast(F)$ is coherent for $i\leq n$.
\end{proposition}

This follows at once from the Leray spectral sequence
\steco
\begin{equation} \label{eq:XII.26} {} \E_2^{p, q}=\R^pf_\ast(\R^qg_\ast(F)) \To \R^\ast h_\ast(F),
\end{equation}
and from the fact that the higher direct images under $f$ of a coherent Module on $X$ are coherent (\EGA III 3.2.1).

\begin{proposition} \label{XII.5.2}
Let $S$ be a locally noetherian prescheme, let $\Ss$ be a quasi-coherent graded Algebra of finite type over $S$, generated by $\Ss_1$, let $X$ be a subprescheme of $\Proj(\Ss)$, let $\OX(1)$ be the invertible Module on $X$ induced by $\Proj(\Ss(1))$ and very ample relative to $S$, let $U$ be an open subset of $X$, let $g:U\to X$ be the canonical immersion, let $h=fg:U\to S$, and let $F$ be a quasi-coherent Module on $U$, giving twisted Modules $F(m)=F\otimes \OX(m)$) ($m\in\ZZ$); let $n$ be an integer and $m_0$ an integer. The following conditions are equivalent:
\begin{enumeratei}
\item
$\R^ig_\ast(F)$ is coherent for $i\leq n$.
\item
$\sisi{\coprod}{\bigoplus}_{m>m_0}\R^ih_\ast(F(m))$ is a finite-type $\Ss$-Module for $i\leq n$.
\end{enumeratei}
\end{proposition}

\begin{proof}
Replacing $F$ by $F(m)$ in the spectral sequence above, we
obtain a spectral sequence of graded $\Ss$-Modules\pageoriginale
\[
\E_2^{p, q}= \sisi{\coprod}{\tbigoplus}_{m\geq m_0} \R^pf_\ast(\R^qg_\ast(F(m)) \To \sisi{\coprod}{\tbigoplus}_{m\geq m_0} \R^\ast h_\ast(F(m)).
\]
Since we have
\[
\R^qg_\ast(F(m)) \simeq \R^qg_\ast(F) (m),
\]
we see that if the $\R^ig_\ast(F)$ are coherent, then $\E_2^{p, q}$ is of finite type over $\Ss$ for $q\leq n$, by part~a) of Lemma~\Ref{XII.5.3} below, which implies that the abutment is of finite type over $\Ss$ in degree $i\leq n$. This proves (i) \ALORS (ii). Moreover, reasoning in the abelian category of graded $\Ss$-Modules modulo the thick subcategory $\mathcal{C}$ of those that are quasi-coherent of finite type, we obtain from the preceding spectral sequence
\[
\sisi{\coprod}{\tbigoplus}_{m\geq m_0} \R^{n+1} h_\ast(F(m)) \simeq \sisi{\coprod}{\tbigoplus}_{m\geq m_0} f_\ast (\R^{n+1} g_\ast(F)(m)) \mod \mathcal{C},
\]
which proves that if the left-hand side is a finite-type $\Ss$-Module, then $\R^{n+1} g_\ast(F)$ is coherent, by part~b) of Lemma~\Ref{XII.5.3}. This proves the implication (ii) \ALORS (i) by induction on $n$. It remains to prove:

\begin{lemme} \label{XII.5.3} Let $S$, $\Ss$, $X$, $f$ be as in \Ref{XII.5.2}, let $G$ be a quasi-coherent Module on $X$, and let $m_0$ be an integer. Then
\begin{enumeratea}
\item
If $G$ is coherent, then for every integer $i$, the graded Module
\[
\sisi{\coprod}{\tbigoplus}_{m\geq m_0} \R^i f_\ast(G(m))
\]
over $\Ss$ is of finite type.
\item
Conversely, suppose that the Module $\sisi{\coprod}{\bigoplus}_{m\geq m_0} \R^i f_\ast(G(m))$ over $\Ss$ is of finite type; then $G$ is coherent.
\end{enumeratea}
\end{lemme}

\begin{proof}[Proof of \Ref{XII.5.3}]
For a), the case $i=0$ is given in \EGA III 2.3.2, and the case $i>0$ in \EGA III 2.2.1 (i)(ii), which says that the $ \R^i f_\ast G(m)$ are coherent\pageoriginale and zero for large $m$ (if we suppose $S$ noetherian, as we may). For b), note that $G$ is isomorphic to $\SheafProj (\sisi{\coprod}{\bigoplus}_{m\geq m_0}f_\ast(G(m))$ (\EGA II 3.4.4 and 3.4.2), which proves that $G$ is coherent if $\sisi{\coprod}{\bigoplus}_{m\geq m_0}f_\ast(G(m))$ is of finite type over $\Ss$, by \loccit 3.4.4.
\skipqed
\end{proof}
\skipqed
\end{proof}

\begin{corollaire}[{{\normalfont$S$ noetherian}}] \label{XII.5.4}
Suppose that $\R^ig_\ast(F)$ is coherent for $i\leq n$. Then, for $i\leq n+1$ and large $m$, there is a canonical isomorphism:
\[
\R^ih_\ast(F(m)) \simeq f_\ast(\R^ig_\ast(G)(m)).
\]
\end{corollaire}

Indeed, the spectral sequence (\Ref{eq:XII.26}) for $F(m)$ then degenerates in degree $\leq n$, by \EGA III 2.2.1~(ii), which gives the result at once (and in fact recovers the implication (ii) \ALORS (i) of~\hyperref[XII.5.2]{5.2}).

\enlargethispage{1.2\baselineskip}%
\begin{corollaire} \label{XII.5.5}
Under the preliminary conditions of \Ref{XII.5.2}, with $S$ noetherian, the following conditions are equivalent:
\begin{enumeratei}
\item
$\sisi{\coprod}{\bigoplus}_{m\geq m_0}h_\ast(F(m)$ is of finite type over $S$, and $\R^ih_\ast(F(m))=0$ for $0<i\leq n$ and large $m$.
\item
$g_\ast(F)$ is coherent, and $\R^ig_\ast(F)=0$ for $0<i\leq n$.
\item[\textup{(ii bis)}] $g_\ast(F)$ is coherent, and $\prof_Y g_\ast(F)>n+1$.
\end{enumeratei}
\end{corollaire}

The equivalence of (ii) and (ii bis) is contained in \Ref{III}~\Ref{III.3.3}. Moreover, by \Ref{XII.5.2}, conditions (i) and (ii) both imply that the $\R^ig_\ast(F)$ ($i\leq n$) are coherent. The equivalence of (i) and (ii) then follows from \Ref{XII.5.4}, taking into account the fact that for every coherent Module $G$ on $X$, one has $G=0$ if and only if $f_\ast(G(m))=0$ for large $m$, for example by \EGA III 2.2.1 (iii).

\begin{remarque} \label{XII.5.6}
The\pageoriginale criteria \Ref{XII.5.2} and \Ref{XII.5.5} can be interpreted by saying that the \og simultaneous finiteness condition\fg \Ref{XII.5.2}~(ii) is expressed by local regularity properties (in terms of depth by \Ref{VIII}~\Ref{VIII.2.1}) of $F$ at the points of $U$ near $Y=X- U$, whereas the \og asymptotic vanishing condition\fg \Ref{XII.5.5}~(i) is distinctly stronger in nature and is expressed by local regularity conditions on $g_\ast(F)$ at the points of~$Y$ itself. It would be interesting, in order to generalize Lefschetz-type theorems for projective morphisms to quasi-projective morphisms, to find local criteria on $X$ that are necessary and sufficient for the $\Ss$-Modules $\sisi{\coprod}{\bigoplus}_{m\geq 0} \R^ih_\ast(F(m))$ for $i\leq n$ to be of finite type. When $S$ is the spectrum of a field (and doubtless more generally if it is the spectrum of an artinian ring) and $Y=X- U$ is finite, one can show that it is necessary and sufficient that the following conditions hold:
\begin{enumerate}
\item[\textup{1\noo)}] $\prof F_x >n$ for every closed point $x$ of $U$ (compare \Ref{XII.1.4}).
\item[\textup{2\noo)}] $\R^ig_\ast(F)$ is coherent for $i\leq n$, or equivalently, there exists an open neighborhood $V$ of $Y$ such that for every closed point $x$ of $U\cap V$, one has $\prof F_x >n+1$.
\end{enumerate}
\end{remarque}

\begin{proposition} \label{XII.5.7}
Let $S$ be a locally Noetherian prescheme, let $g:U\to X$ be a morphism of preschemes of finite type over $S$\sfootnote{It is in fact enough that $g$ be quasi-compact and quasi-separated (\textup{\EGA IV1.2.1}), with no condition on $U,X$.}, with structural morphisms $h$ and $f$, let $F$ be a quasi-coherent module on $U$, and let $n$ be an integer. The following conditions are equivalent:
\begin{enumeratei}
\item
For every base change $S'\to S$, with $S'$ Noetherian, the module $\R^ng'_\ast(F')$ on $X'$ is coherent.
\item
For every base change as above and every coherent ideal $\Jj$ on~$S'$, denoting by $\Ii$ the ideal $\Jj\Oo_{X'}$ on $X'$, the graded module
\[
\sisi{\coprod}{\tbigoplus}_{m\geq 0} \R^ng'_\ast(\Ii^mF')
\]
over $\sisi{\coprod}{\bigoplus}_{m\geq 0} \Ii^m$ is of finite type.
\item
For\pageoriginale every base change $S'\to S$, and $\Jj$ as above, the graded module
\[
\sisi{\coprod}{\tbigoplus}_{m\geq 0} \R^ng'_\ast(\Ii^mF'/\Ii^{m+1}F')
\]
over $\gr_I(\Oo_{X'})=\sisi{\coprod}{\bigoplus}_{m\geq 0} \Ii^m/\Ii^{m+1}$ is of finite type.
\end{enumeratei}
\end{proposition}

Obviously, (ii) \ALORS (i) and (iii) \ALORS (i), as one sees by taking $\Ii=0$ in conditions (ii) and (iii). The reverse implications are obtained by applying (i) to the composite base change $S^{''}\to S'\to S$, where $S^{''}$ is equal to $\Spec \sisi{\coprod}{\bigoplus}_{m\geq 0} \Jj^m$ \resp $\Spec \sisi{\coprod}{\bigoplus}_{m\geq 0} \Jj^m/\Jj^{m+1}$.

The point of this proposition is that conditions of the form (ii) are those that occur in the \og algebraic-formal comparison theorems\fg, whereas conditions of the form (iii) occur in the \og existence theorems\fg that complement them; \cf Expos\'e \Ref{IX}. A first interesting case is that in which $f:X\to S$ is the identity, so that one is dealing with conditions on a morphism $h:U\to S$ locally of finite type and a quasi-coherent module $F$ on $U$, flat over $S$. To obtain sufficient conditions, we shall assume that $U$ embeds by $g:U\to X$ as an open subprescheme of an $X$ that is \emph{proper} over $S$. Applying \Ref{XII.5.1}, we thus see:

\begin{corollaire} \label{XII.5.8}
Let $f:X\to S$ be a proper morphism, with $S$ locally Noetherian, let $U$ be an open subset of $X$, let $g:U\to X$ be the canonical immersion, let $h=fg:U\to S$, and let $F$ be a quasi-coherent module on $X$, flat over $S$. Assume that, for every base change $S'\to S$, with $S'$ locally Noetherian, $\R^ig'_\ast(F')$ is coherent on $X'$ for $i\leq n$. Then the following hold:
\begin{enumeratei}
\item
For every base change $S'\to S$, with $S'$ locally Noetherian, $\R^ih'_\ast(F')$ is coherent on $S'$ for $i\leq n$.
\item
For every $S'\to S$ as above, and every coherent ideal $\Jj$ on $S'$, the graded modules\pageoriginale
\[
\sisi{\coprod}{\tbigoplus}_{m\geq 0} \R^ih'_\ast(\Jj^mF')
\]
over $\sisi{\coprod}{\bigoplus}_{m\geq 0}\Jj^m$ are of finite type for $i\leq n$.
\item
For every $S'\to S$ and $\Jj$ as above, the graded modules
\[
\sisi{\coprod}{\tbigoplus}_{m\geq 0} \R^ih'_\ast(\Jj^mF'/\Jj^{m+1}F')
\]
over $\gr_\Jj(\Oo_{S'})=\sisi{\coprod}{\bigoplus}_{m\geq 0}\Jj^m/\Jj^{m+1}$ are of finite type for $i\leq n$.

Moreover, under the conditions of (ii), and by the comparison theorem \Ref{IX.1.1}, denoting by $\widehat{S'}$ the formal completion of $S'$ with respect to $\Jj$, and by $\widehat{U'}$ that of $U'$ with respect to $\Jj\Oo_{U'}$, the canonical homomorphisms
\[
\widehat{\R^ih'_\ast(F')} \to \R^i\widehat{h'}_\ast(\widehat{F'}) \to \varprojlim_k\R^ih'_\ast(F'_k)
\]
are isomorphisms for $i\leq n-1$.
\end{enumeratei}
\end{corollaire}

\begin{remarque} \label{XII.5.9}
Assume in addition that the conditions of \Ref{XII.5.8} hold with $F$ \emph{coherent}, and consider a base change $S'\to S$ as in \Ref{XII.5.9} (i). Assume further that $S'$ is locally embeddable in a regular scheme, or, more generally, that there exists a faithfully flat and quasi-compact morphism $S^{''}\to S$ such that $S^{''}$ is locally embeddable in a regular scheme; this condition holds in particular if $S'$ is local. Then the conclusion of \Ref{XII.5.8} (i) and (ii) remains valid when $F'$ is replaced there by a module $G'$ on $U'$ such that every point of $U'$ has an open neighborhood on which $G'$ is isomorphic to a module of the form ${F'}^n$. Indeed, one reduces to the case where $S'$ itself is locally embeddable in a regular scheme, and hence so are $S^{''}=\Spec\big(\sisi{\coprod}{\bigoplus}_{m\geq 0}\Jj^m\big)$ and $X^{''}=X'\times_{S'}S^{''}=X\times_SS^{''}$, which are of finite type over it. One then applies the finiteness criterion \Ref{VIII}~\Ref{VIII.2.3} to the direct images for $i\leq n$ of $G^{''}$ under the immersion $U^{''}\to X^{''}$, noting that its conditions hold by hypothesis for $F^{''}$\pageoriginaled, and therefore also for $G^{''}$, since they are expressed in terms of depth and $G^{''}$ is locally isomorphic to an ${F^{''}}^n$. The same argument shows that, if $\Gg'$ is a coherent module on $\widehat{U'}$ (the completion of $U'$ with respect to the ideal $\Jj\Oo_{U'}$), such that $\Gg'_0=\Gg'/\Jj\Gg'$ is locally of the form $F_0^{'n}$, then the conclusion of (iii) remains valid when $F'$ is replaced there by $\Gg'$. Using the results of Expos\'e \Ref{IX}, one thus obtains the following result:
\end{remarque}

\begin{corollaire} \label{XII.5.10}
Let $f:X\to S$ be a proper morphism, with $S$ locally Noetherian, and let $U$ be an open subset of $X$; assume that $U$ is flat over $S$, and that, for every base change $S'\to S$, with $S'$ locally Noetherian, $R^ig'_\ast(\Oo_{U'})$ is coherent on $X'$ for $i=0,1$. Assume now that $S'$ is of the form $\Spec(A)$, where $A$ is a Noetherian ring equipped with an ideal $J$ such that $A$ is separated and complete for the $J$-adic topology. Under these conditions:
\begin{enumeratei}
\item
The functor $F\mto\widehat{F}$ from the category of locally free modules on $U'$ to the category of locally free modules on $\widehat{U'}$ is fully faithful.
\item
For every locally free module $\formelF$ on $\widehat{U'}$, there exists a coherent module $F$ on~$U'$ (not necessarily locally free, alas), and an isomorphism $\widehat{F}\simeq\formelF$.
\end{enumeratei}
\end{corollaire}

It remains only to prove (ii), by \Ref{XII.5.9}. But it follows from this remark and \Ref{IX.2.1} that $\formelF$ is induced by a coherent module $\formelG$ on $\widehat{X'}$. By the existence theorem \EGA III 5.1.4, $\formelG$ is of the form $\widehat{F}$, where $F$ is coherent on $X$, which proves the conclusion.

\begin{remarques} \label{XII.5.11}
1\textsuperscript{o}) Using \Ref{XII.5.10}, \Ref{XII.4.7}, and a suitable hypothesis saying that certain local rings of the geometric fibers of $X'\to S'$ are \og pure\fg \resp parafactorial, one should be able to obtain statements saying that the functor $Z'\mto\widehat{Z'}$ from the category of finite \`etale coverings of $X'$ to the category of finite \`etale coverings\pageoriginale of $\widehat{X'}$ (or, equivalently, of $X'_0$) is an equivalence of categories, \resp that the functor $L\mto\widehat{L}$ from the category of invertible modules on $X'$ to the category of invertible modules on $\widehat{X'}$ is an equivalence. Using recent results of Murre, it is likely that one should be able to deduce existence theorems for Picard schemes of certain \emph{nonproper} algebraic schemes\nde{Of course, in the projective case one should refer to Grothendieck's existence theorems in FGA; \cf Grothendieck~A., {\og Technique de descente et th\'eor\`emes d'existence en g\'eom\'etrie alg\'ebrique. VI. Les sch\'emas de Picard: propri\'et\'es g\'en\'erales\fg}, in \emph{S\'eminaire Bourbaki}, vol.~7, Soci\'et\'e math\'ematique de France, Paris, 1995, \Exp 236, p\ptbl 221--243 and {\og Technique de descente et th\'eor\`emes d'existence en g\'eom\'etrie alg\'ebrique. V. Les sch\'emas de Picard: th\'eor\`emes d'existence\fg}, in \emph{S\'eminaire Bourbaki}, vol.~7, Soci\'et\'e math\'ematique de France, Paris, 1995, \Exp 232, 143--161. The nine finiteness conjectures found there are proved in Expos\'es \Ref{XII} and \Ref{XIII} by Mme Raynaud and Kleiman in \SGA 6. For an excellent elementary account of the subject, see Kleiman's expository article (Kleiman~S., {\og The Picard scheme\fg}, forthcoming in Contemp. Math.). For an application of these techniques to the global generalized Jacobians of a relative smooth curve, see (Contou-Carr\`ere~C., {\og La jacobienne g\'en\'eralis\'ee d'une courbe relative; construction et propri\'et\'e universelle de factorisation\fg}, \textit{C.~R. Acad. Sci. Paris S\'er.~A-B} \textbf{289} (1979), \numero3, A203--A206 and {\og Jacobiennes g\'en\'eralis\'ees globales relatives\fg}, in \emph{The Grothendieck Festschrift, Vol. II}, Progr. Math., vol.~87, Birkh\"auser, Boston, 1990, p\ptbl 69--109). See also, by the same author and in the purely local setting, the construction and study of the \og local generalized Jacobian\fg functor ({\og Jacobienne locale, groupe de bivecteurs de Witt universel, et symbole mod\'er\'e\fg}, \textit{C.~R. Acad. Sci. Paris S\'er.~I Math.} \textbf{318} (1994), \numero 8, p\ptbl 743--746). Moreover, although in the case of a projective smooth morphism the connected components of the Picard scheme are proper, this is no longer so in the singular case. The problem of compactifying Picard schemes then arises naturally; this problem has been studied in detail, notably in (Altman~A.B. \& Kleiman~S., {\og Compactifying the Picard scheme\fg}, \emph{Adv. in Math.} \textbf{35} (1980), \numero 1, p\ptbl 50--112, and {\og Compactifying the Picard scheme. II\fg}, \emph{Amer. J.~Math.} \textbf{101} (1979), \numero 1, p\ptbl 10--41). The case of curves had been studied earlier (D'Souza~C., {\og Compactification of generalised Jacobians\fg}, \emph{Proc. Indian Acad. Sci. Sect. A Math. Sci.} \textbf{88} (1979), \numero 5, p\ptbl 419--457). One even knows exactly when the compactified Jacobian of a curve is irreducible (Rego~C.J., {\og The compactified Jacobian\fg}, \emph{Ann. Sci. \'Ec. Norm. Sup. (4)} \textbf{13} (1980), \numero 2, p\ptbl 211--223), the latter being the closure of the ordinary Jacobian when the curve is geometrically integral and locally planar; for a construction in families of compactified Jacobians, see (Esteves~E., {\og Compactifying the relative Jacobian over families of reduced curves\fg}, \emph{Trans. Amer. Math. Soc.} \textbf{353} (2001), \numero 8, p\ptbl 3045--3095). Since then, existence results for the Picard scheme in the proper case have advanced beyond the original edition of \SGA 2; \cf (Murre~J.P., {\og On contravariant functors from the category of pre-schemes over a field into the category of abelian groups (with an application to the Picard functor)\fg}, \emph{Publ. Math. Inst. Hautes \'Etudes Sci.} \textbf{23} (1964), p\ptbl 5--43), and especially (Artin~M., {\og Algebraization of formal moduli. I\fg}, in \emph{Global Analysis (Papers in Honor of K\ptbl Kodaira)}, Univ. Tokyo Press, Tokyo, 1969, p\ptbl 21--71). See also (Raynaud~M., {\og Sp\'ecialisation du foncteur de Picard\fg}, \emph{Publ. Math. Inst. Hautes \'Etudes Sci.} \textbf{38} (1970), p\ptbl 27--76) for a scheme proper over a discrete valuation ring but not necessarily flat. For a discussion of more recent results, in particular Artin's, for the Picard functor of proper flat schemes, especially in the cohomologically flat case in dimension $0$, see Chapter VIII of (Bosch~S., L\"utkebohmert~W. \& Raynaud~M., \emph{N\'eron models}, Ergebnisse der Mathematik und ihrer Grenzgebiete (3), vol.~21, Springer-Verlag, Berlin, 1990) and the references cited there. Much more recently, very precise results have been obtained for relative curves $f:X\to S$ \sisi{over}{over} the spectrum $S$ of a discrete valuation ring with perfect residue field. More precisely, one assumes that $f$ is proper and flat, that $X$ is regular, and that $f_*\OX=\Oo_S$. On the other hand, one does not assume that $f$ is cohomologically flat in dimension~$0$, \ie one does not assume that $H^1(X,\Oo)$ is torsion-free. The Picard functor is then not representable, either by a scheme or by an algebraic space, as soon as the torsion in question is nonzero. Let~$J$ be the N\'eron model of the generic fiber of $f$; it is a quotient of the Picard functor $P$. Raynaud then showed that the kernel of the tangent map $H^1(X,\Oo)=\mathrm{Lie}(P)\to\mathrm{Lie}(J)$ coincides with the torsion subgroup of $H^1$, and that the cokernel has the same length (see Theorem 3.1 of (Liu~Q., Lorenzini~D. \& Raynaud~M., {\og N\'eron models, Lie algebras, and reduction of curves of genus one\fg}, \emph{Invent. Math.} \textbf{157} (2004), p\ptbl 455--518)). This result enables those authors to study the relation between the Birch--Swinnerton-Dyer and Artin--Tate conjectures (see Th.\ptbl 6.6 of {\loccit}). As for the local Picard scheme, see Boutot's thesis, cited in editor's note~\eqref{Boutot} on page~\pageref{Boutot}.}. In general, the removal of purity hypotheses in various existence theorems, especially representability theorems for functors such as the Hilbert or Picard functors, by means of the techniques developed in this seminar, deserves systematic study.

2\textsuperscript{o}) One may seek manageable necessary and sufficient conditions, in terms of depth, for the universal finiteness condition considered in \Ref{XII.5.10} to hold. When $S$ is the spectrum of a field, it follows easily from \EGA III 1.4.15 that it is necessary and sufficient for the $\R^ig_\ast(F)$ ($i\leq n$) to be coherent, which is indeed expressed in terms of depth by \Ref{VIII}~\Ref{VIII.2.3}. In the general case, however, note that it is not enough to require the preceding condition to hold for all the fibers $U_s\subset X_s$ ($s\in S$), even when $n=0$. For example, take $X=S$, let $S$ be the spectrum of a discrete valuation ring, let $U$ be the open subset consisting only of the generic point, and let $F=\Oo_U$.

3\textsuperscript{o}) Here, however, is a \emph{sufficient} condition ensuring that the hypotheses of \Ref{XII.5.10} hold: it is enough that $f$ be flat and that, for every $s\in S$ and every $x\in Y_s=X_s-U_s$, one have
\[
\prof \Oo_{X_s,x}\geq n+2.
\]
Indeed, taking Lemma \Ref{XII.2.5} into account (\cf relation (\Ref{eq:XII.16}) following \Ref{XII.2.5}), it follows that $g_*(\Oo_U)\simeq\OX$ and $\R^ig_*(\Oo_U)=0$ for $0<i\leq n$, and the same relations will clearly hold after every base change $S'\to S$.
\end{remarques}
